\chapter{Metodología}
\label{cap:metodologia}

\section{Enfoque y diseño}
\label{sec:met-enfoque}

La investigación es cuantitativa, experimental y aplicada. Es experimental porque manipula de forma deliberada el estado de frecuencia de CPU y GPU sobre un conjunto controlado de cargas y observa su efecto sobre el tiempo, la energía y el Producto Energía--Retardo. Es aplicada porque su producto es un agente de software funcional cuya utilidad se evalúa empíricamente frente a gobernadores nativos del sistema operativo y del controlador de la GPU. El trabajo se organiza en cuatro fases encadenadas (Figura \ref{fig:met-fases}): la Fase 1 produce el conjunto de datos etiquetado; la Fase 2 entrena y valida los clasificadores y deriva las tablas de política; la Fase 3 implementa el agente; la Fase 4 lo evalúa. Los esquemas y cuadros de verificación complementarios se reúnen en el Anexo \ref{anexo:metodologia}.

\begin{figure}[H]
\centering
\begin{tikzpicture}[
    font=\footnotesize,
    fase/.style={draw, rounded corners=2pt, align=center, minimum height=1.5cm, text width=3.1cm, inner sep=4pt, fill=blue!7, draw=blue!45!black},
    dato/.style={draw, rounded corners=2pt, align=center, minimum height=0.9cm, text width=3.1cm, inner sep=3pt, fill=orange!10, draw=orange!55!black},
    flujo/.style={-{Latex[length=2.2mm]}, semithick},
]
    \node[fase] (f1) {\textbf{Fase 1}\\Caracterización\\y datos etiquetados};
    \node[fase, right=1.0cm of f1] (f2) {\textbf{Fase 2}\\Clasificadores\\y tablas de política};
    \node[fase, right=1.0cm of f2] (f3) {\textbf{Fase 3}\\Agente de control\\DVFS};
    \node[fase, right=1.0cm of f3] (f4) {\textbf{Fase 4}\\Evaluación\\frente a REF};
    \node[dato, below=1.0cm of f1] (d1) {Conjunto de datos\\CPU y GPU\\con etiqueta Roofline};
    \node[dato, below=1.0cm of f2] (d2) {Modelos ONNX\\y política de\\frecuencia};
    \node[dato, below=1.0cm of f3] (d3) {Procesos de CPU\\y de GPU con\\tres brazos};
    \node[dato, below=1.0cm of f4] (d4) {EDP, energía\\y exactitud de\\clasificación};
    \foreach \a/\b in {f1/f2,f2/f3,f3/f4} \draw[flujo] (\a) -- (\b);
    \foreach \a/\b in {f1/d1,f2/d2,f3/d3,f4/d4} \draw[flujo] (\a) -- (\b);
\end{tikzpicture}
\caption[Fases de la metodología]{Fases de la metodología y producto de cada una. Fuente: elaboración propia.}
\label{fig:met-fases}
\end{figure}

\subsection{Población, muestra y criterios de inclusión}

La población son las aplicaciones para sistemas heterogéneos CPU--GPU; se estudió una muestra intencional de cargas reproducibles con regímenes contrastantes. Cada carga debía verificar su resultado, exponer trabajo de punto flotante medible por el instrumento y admitir una etiqueta de intensidad operacional. La unidad de diversidad y de validación fue la familia algorítmica: variantes de tamaño del mismo algoritmo no se trataron como familias independientes. El catálogo y su cribado se detallan en el Anexo \ref{anexo:metodologia}.

\subsection{Plataforma experimental}

Todos los experimentos se ejecutan en un único nodo, con asignación exclusiva del gestor de recursos para que ninguna carga ajena comparta el procesador durante las mediciones (Tabla \ref{tab:met-plataforma}). La plataforma se eligió porque expone la energía por interfaz interna sin privilegios administrativos (RAPL) y un acelerador instrumentable con NVML. Las actuaciones de frecuencia de CPU se restringen a los procesadores lógicos delegados, los núcleos 0 a 5 de un zócalo y sus hilos hermanos, porque el dominio de frecuencia es el núcleo físico y escribir un solo hilo no fija el reloj (anexo, sección \ref{sec:dominios}).

\begin{table}[H]
\centering
\caption{Plataforma experimental.}
\label{tab:met-plataforma}
\small
\begin{tabular}{@{}>{\raggedright\arraybackslash}p{4.2cm}>{\raggedright\arraybackslash}p{9.6cm}@{}}
\toprule
\textbf{Componente} & \textbf{Descripción} \\
\midrule
CPU & 2 $\times$ Intel Xeon Gold 5315Y (Ice Lake-SP), 8 núcleos físicos por zócalo, 2 hilos por núcleo, 2 nodos NUMA \\
Frecuencia & \texttt{intel\_pstate}, 0.8 a 3.2~GHz (3.6~GHz con \textit{turbo} activo), dominio por núcleo físico \\
Energía de CPU & Intel RAPL, dominios de paquete y de DRAM por zócalo \\
GPU & NVIDIA A100 PCIe 40~GB, telemetría por NVML \\
Ejecución & Nodo en asignación exclusiva (Slurm) \\
\bottomrule
\end{tabular}
\end{table}

\subsection{Instrumento de medición}

El instrumento propio tiene dos capas (Figura \ref{fig:arquitectura-instrumento}). La capa inferior, en C++, lanza la carga como proceso hijo detenido antes de su primera instrucción, abre los contadores de hardware sobre ese proceso con herencia y solo entonces lo reanuda, de modo que ninguna instrucción queda fuera de la observación y la medida abarca todo el árbol de procesos e hilos de la carga. Un hilo productor lee los contadores a intervalos regulares y deposita las muestras en una estructura circular de un productor y un consumidor; un hilo consumidor, en otro procesador lógico, las persiste (Figura \ref{fig:muestreo-spsc}, anexo \ref{anexo:metodologia}). La capa superior, en Python, ejecuta la lógica experimental: interpreta el manifiesto de la campaña, verifica la plataforma, calibra, planifica cada combinación, aplica el estado de frecuencia, valida cada corrida y construye el conjunto de datos.

\begin{figure}[tb]
\centering
\resizebox{0.95\textwidth}{!}{
\begin{tikzpicture}[
    font=\small,
    st/.style={draw, rounded corners=1.5pt, align=center,
        minimum height=0.95cm, text width=2.35cm, inner sep=2pt},
    py/.style={st, fill=blue!7, draw=blue!45!black},
    cpp/.style={st, fill=orange!10, draw=orange!55!black},
    ar/.style={-{Latex[length=1.8mm]}, semithick},
    cr/.style={-{Latex[length=1.8mm]}, semithick, dashed},
    sub/.style={font=\tiny, text=black!70, align=center},
    zlab/.style={font=\tiny\bfseries, text=black!55},
]
    % ---------- Banda A: capa Python, preparación (izq -> der) ----------
    \node[py] (man)  {Manifiesto de\\campaña\\{\tiny manifest}};
    \node[py, right=0.35cm of man]  (pre)  {Verificaciones\\previas\\{\tiny preflight}};
    \node[py, right=0.35cm of pre]  (cal)  {Calibración\\Roofline\\{\tiny calibration}};
    \node[py, right=0.35cm of cal]  (plan) {Matriz, orden\\de ejecución\\{\tiny campaign}};
    \node[py, right=0.35cm of plan] (loop) {\textbf{Bucle por}\\\textbf{combinación}\\{\tiny runner, freqctl}};
    \draw[ar] (man) -- (pre);
    \draw[ar] (pre) -- (cal);
    \draw[ar] (cal) -- (plan);
    \draw[ar] (plan) -- (loop);

    % ---------- Banda B: capa C++, medición (der -> izq) ----------
    \node[cpp, below=1.7cm of loop] (exec)
        {Ejecutor de la carga\\{\tiny hijo detenido,}\\{\tiny adjuntado y reanudado}\\{\tiny kernel\_launcher}};
    \node[cpp, left=0.35cm of exec] (coll)
        {Colector y lectores\\{\tiny perf, uncore,}\\{\tiny RAPL, NVML}};
    \node[cpp, left=0.35cm of coll] (ring)
        {Estructura circular\\{\tiny (Figura \ref{fig:muestreo-spsc})}\\{\tiny spsc\_ring}};
    \draw[ar] (exec) -- (coll);
    \draw[ar] (coll) -- (ring);
    \draw[cr] (loop.south) -- node[left, sub]
        {invoca por subproceso;\\antes fija la frecuencia} (exec.north);

    % ---------- Banda C: capa Python, consolidación (izq -> der) ----------
    \node[py, below=1.9cm of ring] (val) {Validación de\\la corrida\\{\tiny validation}};
    \node[py, right=0.35cm of val]  (post) {Post-proceso\\{\tiny postprocess}};
    \node[py, right=0.35cm of post]  (data)
        {\textbf{Conjunto de datos}\\{\tiny windows, intervalos}\\{\tiny \textit{uncore}, fases GPU}};
    \draw[cr] (ring.south) -- node[right, sub]
        {muestras crudas\\{\tiny (samples.csv)}} (val.north);
    \draw[ar] (val) -- (post);
    \draw[ar] (post) -- (data);

    % ---------- zonas ----------
    \begin{pgfonlayer}{background}
        \node[draw=blue!35, fill=blue!3, rounded corners, inner sep=0.2cm,
            fit=(man)(loop), label={[zlab, blue!45!black]above left:Capa Python: preparación y planificación}] {};
        \node[draw=orange!45, fill=orange!4, rounded corners, inner sep=0.2cm,
            fit=(exec)(ring), label={[zlab, orange!55!black]above left:Capa C++: medición}] {};
        \node[draw=blue!35, fill=blue!3, rounded corners, inner sep=0.2cm,
            fit=(val)(data), label={[zlab, blue!45!black]below left:Capa Python: validación y consolidación}] {};
    \end{pgfonlayer}
\end{tikzpicture}
}

\caption[Recorrido de una campaña a través de las dos capas del instrumento]{Recorrido de una
campaña a través de las dos capas del instrumento. La capa Python prepara la campaña (banda
superior), invoca por subproceso el ejecutor de la capa C++ una vez por combinación de carga y
frecuencia (banda intermedia), y consolida las muestras crudas que este produce en el conjunto de
datos (banda inferior). La capa C++ concentra la medición sensible a la latencia; la capa Python,
la lógica experimental, donde priman la claridad y la trazabilidad. Fuente: elaboración propia.}
\label{fig:arquitectura-instrumento}
\end{figure}

La Tabla \ref{tab:met-senales} resume las señales. Los contadores de CPU se limitan a nueve eventos simultáneos, dimensionados contra el presupuesto de contadores físicos del nodo, verificado empíricamente, para evitar la multiplexación y su extrapolación. Una \textbf{ventana de muestreo} es el intervalo entre dos lecturas consecutivas; su contenido es la diferencia de las lecturas acumuladas. El tráfico a DRAM se mide con los contadores del controlador de memoria (\textit{uncore}), cuyo intervalo práctico es de 10~ms o más, ya que el manual de perf advierte de overhead crítico y desestabilidad sub-10~ms \cite{PerfMeasurementOverhead2026}. En GPU, el recolector solicita lecturas NVML cada 5~ms mediante \texttt{gpu\_interval\_ns=5\,000\,000}; sin embargo, la actualización física de los sensores del controlador tiene una cadencia efectiva aproximada de 100--120~ms. Por ello, las muestras NVML se usan como telemetría agregada y no como ventanas GPU con resolución de 5~ms.

\subsection{Modelo Roofline y etiquetado}
\label{sec:met-roofline}

La etiqueta de cada observación (\textit{compute-bound} o \textit{memory-bound}) resulta de comparar su intensidad operacional $I$ con el punto de inflexión del modelo Roofline (ecuación \eqref{eq:ridge}), calibrado en la plataforma por precisión aritmética y, en las campañas con barrido, por nivel de frecuencia. Los techos se calibran al inicio de cada campaña, con la misma asignación de núcleos que las corridas, mediante microbenchmarks de ancho de banda y de cómputo, y se contrastan con una referencia declarada. La etiqueta se deriva sin juicio manual, y la expectativa cualitativa de régimen que declara cada carga se usa solo como contraste.

Los dos dispositivos obtienen $I$ de forma distinta (Figura \ref{fig:met-etiqueta}). En CPU, los FLOPs se miden por ventana con los sub-eventos de \texttt{FP\_ARITH\_INST\_RETIRED} ponderados por el ancho vectorial (ecuación \eqref{eq:flops-medidos}), y los bytes con los comandos CAS del controlador de memoria; $I$ se calcula por intervalo, en vivo. En las campañas de GPU empleadas en este trabajo, NVML no proporciona una medición de FLOPs ni de bytes por ventana o por lanzamiento. Por ello, la intensidad operacional se caracteriza fuera de línea con Nsight Compute \cite{Nsight2026} (ecuación \eqref{eq:intensidad-gpu}) y se asocia a cada corrida. Una etiqueta con granularidad de fase requeriría combinar el perfilado de contadores con una traza temporal de lanzamientos mediante CUPTI, mecanismo que no forma parte de las observaciones aceptadas para este conjunto. Para un kernel determinista, con tamaño, precisión, entradas y configuración de lanzamiento fijos, esta caracterización se trata como una propiedad operacional estable tras verificar la convergencia entre los dos puntos de mayor trabajo, cuyo cambio relativo debe ser inferior al 1\%. Esta condición no afirma que el tráfico observado sea invariante ante cualquier configuración de ejecución. La caracterización distingue la precisión efectivamente ejercitada y excluye las cargas de aritmética entera o de precisión mixta no resoluble.

\begin{figure}[H]
\centering
\begin{tikzpicture}[
    font=\footnotesize,
    b/.style={draw, rounded corners=2pt, align=center, minimum height=1.0cm, text width=3.3cm, inner sep=3pt},
    cpu/.style={b, fill=blue!7, draw=blue!45!black},
    gpu/.style={b, fill=orange!10, draw=orange!55!black},
    et/.style={b, fill=green!10, draw=green!45!black},
    fl/.style={-{Latex[length=2.2mm]}, semithick},
    lb/.style={font=\scriptsize\bfseries, text=black!60},
]
    \node[cpu] (c1) at (0,1.5) {FLOPs por ventana\\(PMU)};
    \node[cpu] (c2) at (0,-0.1) {Bytes por intervalo\\(CAS del IMC)};
    \node[cpu] (c3) at (4.6,0.5) {$I$ por intervalo\\(en vivo)};
    \node[gpu] (g1) at (0,-2.0) {Nsight Compute\\(fuera de línea)};
    \node[gpu] (g2) at (4.6,-2.0) {$I$ fija por kernel\\(tras convergencia)};
    \node[et] (r) at (9.2,-0.75) {$I$ frente al punto\\de inflexión: etiqueta};
    \node[lb, anchor=south] at (2.3,1.85) {CPU};
    \node[lb, anchor=north] at (2.3,-2.7) {GPU};
    \draw[fl] (c1.east) -- (c3.west);
    \draw[fl] (c2.east) -- (c3.west);
    \draw[fl] (g1) -- (g2);
    \draw[fl] (c3.east) -- (r.north west);
    \draw[fl] (g2.east) -- (r.south west);
\end{tikzpicture}
\caption[Obtención de la etiqueta en CPU y en GPU]{Obtención de la etiqueta Roofline. En CPU la intensidad se mide en vivo por intervalo del controlador de memoria; en GPU se caracteriza una vez por kernel, fuera de línea. Fuente: elaboración propia.}
\label{fig:met-etiqueta}
\end{figure}

\section{Fase 1: recolección y caracterización}
\label{sec:met-fase1}

\subsection{Verificación previa y catálogo}

Antes de medir, una batería de verificaciones automáticas, parte del comando que lanza la campaña, confirma las capacidades del entorno, la integridad del catálogo y la disponibilidad de la instrumentación (Tabla \ref{tab:verificaciones-previas}, anexo \ref{anexo:metodologia}). Las verificaciones bloqueantes detienen la campaña antes de tocar el sistema: se prefiere no producir datos a producir datos cuya validez no puede sostenerse. El catálogo definitivo es una lista versionada y congelada de familias elegibles, tras un cribado con una rejilla de frecuencia reducida de tres niveles que descarta candidatos, comprueba la utilidad de su señal y estima su cobertura del plano Roofline. Cada clase debe quedar representada por al menos cinco o seis familias distintas por dispositivo. Con esto en mente, el conjunto de CPU reúne 30 familias, y el de GPU 16 (anexo, secciones \ref{sec:banco-candidatos} y \ref{sec:cribado}).

\subsection{Matriz experimental y control de sesgos}

La matriz resulta de combinar las cargas del catálogo, los estados de frecuencia (Tabla \ref{tab:frecuencias} del Anexo \ref{anexo:metodologia}) y las repeticiones independientes del binario. Los niveles se verifican contra los relojes soportados por el nodo. En CPU, fijar un nivel requiere acotar el rango por núcleo y desactivar el \textit{turbo} global para evitar ajustes autónomos del gestor de energía. Cada ventana se clasifica como válida, no confiable, no verificable (margen tras transición) o no aplicable (nivel nativo).

Para mitigar la latencia de asentamiento del reloj, la cual fue medida bajo carga en 10–11~ms para CPU y 70–80~ms para GPU, se verificó cada nivel dentro de una tolerancia de $\pm5\%$ y se aplicó un margen de guarda de 0.15~s al inicio y final de la región medible.

Las medidas contra el sesgo son cinco. Cada combinación se repite tres veces (validado empíricamente como suficiente). El orden de ejecución se aleatoriza con una semilla registrada en el manifiesto. Se incluyen repeticiones de una carga de referencia, cuya dispersión relativa señala inestabilidad del entorno. El eje de frecuencia de un dispositivo se mide con el otro dispositivo en su \textit{governor} nativo, porque fijar la CPU en su nivel mínimo mientras se barre la GPU alarga las cargas de GPU y contamina la comparación. Y toda campaña reanudable conserva una huella del protocolo, de modo que dos manifiestos con protocolos distintos no se mezclan (anexo, sección \ref{sec:matriz-sesgos}).


\subsection{Calidad de los datos y granularidad}

Ninguna observación se descarta en silencio. Cada ventana conserva una anotación de calidad (válida, sin diferencia previa, excluida por calentamiento, contadores degradados, intensidad indefinida, energía inválida o sin lectura de frecuencia) y el filtrado es una decisión del análisis. El transitorio de arranque de cada carga se identifica en dos etapas sobre una señal representativa del estado de carga (IPC en CPU, utilización en GPU): el primer instante desde el cual el coeficiente de variación se mantiene en 5\% o menos en dos ventanas móviles consecutivas y, como respaldo, un método de detección de puntos de cambio. Ese intervalo solo actúa como filtro posterior y nunca condiciona la adquisición. Cada corrida completa se valida además por integridad global: sin pérdidas en la estructura de muestreo, suma de verificación correcta y criterio de éxito propio cumplido (anexo, sección \ref{sec:calentamiento}).

La unidad de observación del entrenamiento se declara por dispositivo (Tabla \ref{tab:met-granularidad}), porque en ambos casos la ventana del instrumento no coincide con la unidad que posee una etiqueta con respaldo físico. En CPU es un \textbf{intervalo del controlador de memoria}, con las variables recalculadas sobre los deltas agregados del intervalo y no promediando cocientes. En GPU es una \textbf{corrida completa}, con agregados robustos (mediana, dispersión) de las muestras válidas posteriores al calentamiento. Las lecturas repetidas de una misma corrida no se tratan como ejemplos independientes.

\subsection{Reproducibilidad}

Cada campaña queda definida por un manifiesto declarativo bajo control de versiones (cargas, estados de frecuencia, repeticiones, período de muestreo, asignación de núcleos, semilla y referencias de calibración). Junto con los datos se persisten el perfil de capacidades del nodo, la calibración y los metadatos de cada corrida. Una campaña interrumpida se reanuda saltando las corridas ya aceptadas, protegida por la huella del protocolo (anexo, sección \ref{sec:reproducibilidad}).

\section{Fase 2: modelos de clasificación y política de frecuencia}
\label{sec:met-fase2}

\subsection{Formulación, variables y validación}

Se formula una clasificación binaria supervisada por dispositivo; la salida es la probabilidad de la clase \textit{memory-bound}. En CPU, las seis variables son cinco tasas derivadas de los contadores del intervalo (IPC, MPKI, tasa de fallos de caché, fracción de ciclos detenidos por memoria e instrucciones por segundo) y la frecuencia observada, conservada como covariable porque la misma carga puede cambiar de posición respecto del punto de inflexión cuando cambia el techo de cómputo. Antes de entrenar, se revisan las correlaciones entre las columnas candidatas. Los pares con correlación absoluta superior a 0.85 (común en la literatura), y la colinealidad restante, evaluada mediante el factor de inflación de la varianza, orientan el descarte de variables redundantes. El conjunto final se registra en un contrato versionado por dispositivo (anexo, sección \ref{sec:contrato-features}).

La validación es \textit{leave-one-familia-out}: en cada pliegue se retira por completo una familia, que no participa en el ajuste, la selección de representación ni la elección del umbral (Figura \ref{fig:met-lofo}). Las celdas de familia y clase se limitan a 1000 intervalos con una semilla fija, y cada observación se pondera de forma inversa al tamaño de su celda (peso medio uno), de modo que ninguna celda domina por volumen. La métrica principal es la exactitud balanceada por celda de familia y clase,
\begin{equation}
\mathrm{EB} = \frac{1}{|\mathcal{C}|}\sum_{(f,c)\in\mathcal{C}} \frac{a_{fc}}{n_{fc}},
\label{eq:met-eb}
\end{equation}
donde $\mathcal{C}$ es el conjunto de celdas observadas (cada combinación de familia $f$ y clase $c$), $|\mathcal{C}|$ es el número de celdas, $a_{fc}$ son los aciertos en cada celda y $n_{fc}$ las observaciones totales de esa celda. Está definida aunque una familia tenga una sola clase y no depende de la prevalencia. La incertidumbre se reporta por dos vías: cinco semillas de submuestreo, y remuestreo con reemplazo de las familias (2000 réplicas) para los intervalos de confianza del 95\% y las comparaciones pareadas.

\begin{figure}[H]
\centering
\begin{tikzpicture}[
    font=\footnotesize,
    b/.style={draw, rounded corners=2pt, align=center, minimum height=0.9cm, inner sep=3pt},
    ext/.style={b, fill=blue!7, draw=blue!45!black, text width=4.2cm},
    mid/.style={b, fill=orange!10, draw=orange!55!black, text width=4.2cm},
    sal/.style={b, fill=green!10, draw=green!45!black, text width=4.2cm},
    fl/.style={-{Latex[length=2.2mm]}, semithick},
]
    \node[ext] (a) {Pliegue externo:\\se retira la familia $f$};
    \node[mid, right=0.5cm of a] (b) {LOFO interno sobre las\\demás familias: elige\\el umbral $\tau$};
    \node[ext, right=0.5cm of b] (c) {Se ajusta el modelo\\en las demás familias};
    \node[sal, below=0.8cm of a] (e) {EB, cobertura\\y abstenciones};
    \node[sal, below=0.8cm of c] (d) {Se predice $f$ con\\la regla selectiva};
    \draw[fl] (a) -- (b); \draw[fl] (b) -- (c); \draw[fl] (c) -- (d); \draw[fl] (d) -- (e);
\end{tikzpicture}
\caption[Validación anidada por familia]{Validación anidada \textit{leave-one-familia-out}. La familia retirada no interviene en ninguna decisión de ajuste. Fuente: elaboración propia.}
\label{fig:met-lofo}
\end{figure}

\subsection{Decisión selectiva y calibración}

Con $p_i$ la probabilidad de \textit{memory-bound} y $q_i=\max(p_i,1-p_i)$ la confianza, la regla selectiva entrega la clase solo si $q_i\geq\tau$ y en caso contrario emite \textit{revisar}:
\begin{equation}
\widehat{y}_i=\begin{cases}
\textit{memory-bound}, & p_i\geq 0.5\ \land\ q_i\geq\tau,\\
\textit{compute-bound}, & p_i<0.5\ \land\ q_i\geq\tau,\\
\textit{revisar}, & q_i<\tau.
\end{cases}
\label{eq:met-selectiva}
\end{equation}
El umbral $\tau$ se elige sobre una rejilla mediante un LOFO interno, maximizando la exactitud balanceada por celda calculada únicamente sobre las observaciones aceptadas, con una cobertura media por celda mínima de 0.60. Las salidas \textit{revisar} se consideran abstenciones y se excluyen del cálculo de las métricas de clasificación sobre las predicciones aceptadas. Su efecto se cuantifica mediante la cobertura, definida como la proporción del total de observaciones para las cuales el modelo emite una clase. La confianza se evalúa además mediante el error de calibración esperado, que divide los valores de confianza en diez rangos de igual amplitud y compara, en cada rango, la confianza media con la proporción real de aciertos,
\begin{equation}
\mathrm{ECE} = \sum_{b=1}^{10} \frac{n_b}{N}\left|\mathrm{conf}_b - \mathrm{acc}_b\right|,
\label{eq:met-ece}
\end{equation}
donde $N$ es el número total de observaciones evaluadas, $n_b$ es el número de observaciones del rango $b$, $\mathrm{conf}_b$ es la confianza media de ese rango y $\mathrm{acc}_b$ es su proporción de aciertos.

\subsection{Comparación y selección de modelos}

Se compararon siete modelos bajo el mismo protocolo: una regla mayoritaria, un árbol de profundidad uno, uno de profundidad seis, regresión logística, \textit{random forest}, \textit{extra trees} y XGBoost, con hiperparámetros fijos. La búsqueda de hiperparámetros con Optuna se evaluó anidada en cada pliegue externo y se descartó: con un número limitado de familias, el criterio interno que maximiza no discrimina de forma confiable la configuración que generaliza mejor. Para explorar si el desempeño observado estaba limitado principalmente por la capacidad del modelo o por la información contenida en las variables se aplicaron cuatro diagnósticos complementarios. Entre ellos se evaluó un clasificador de vecinos más cercanos sin ajuste, utilizado para examinar la separabilidad local de las clases en el espacio de variables (anexo, sección \ref{sec:metodologia-limite-cpu}). La selección del modelo final se realizó después de comparar el desempeño bajo el protocolo LOFO y de verificar la invarianza de las variables frente a la acción del agente. Los modelos y umbrales seleccionados se reportan en el capítulo de resultados.

\subsection{Invarianza de las entradas frente a la acción del agente}
\label{sec:metodologia-invarianza-gpu}

La comparación por exactitud identifica candidatos que generalizan fuera de familia, pero un clasificador integrado en un controlador debe usar entradas que no cambien como consecuencia de su propia acción. En GPU se audita esa condición con dos pruebas sobre los datos ya disponibles. La prueba contrafactual sustituye el valor de una entrada potencialmente afectada por la actuación y cuantifica los cambios de predicción. La prueba de exclusión repite la validación LOFO sin dicha entrada. Toda variable que no supera ambas pruebas se excluye del contrato operativo.

Los modelos se vuelven a evaluar con el contrato de entradas resultante. La selección exige exactitud balanceada, calibración y cobertura suficientes bajo LOFO; la latencia de inferencia se informa como propiedad de implementación, mientras que la latencia completa desde la detección hasta la actuación se mide separadamente durante la evaluación del agente. En CPU se aplica la misma auditoría a la frecuencia observada, mediante una sustitución contrafactual coherente con el cambio de frecuencia que puede ejecutar el agente. Los resultados de ambas auditorías se presentan en las secciones \ref{sec:resultados-gpu-invarianza} y \ref{sec:resultados-cpu-frecuencia-contrafactual}.

\subsection{Derivación de la política de frecuencia}

La política de frecuencia se deriva de la campaña de la Fase 1, por separado del clasificador. Cada nivel se compara con REF sobre el mismo kernel y la ganancia agregada se estima por \textit{bootstrap} de kernels. La regla de elección se valida dejando una familia fuera; cualquier refinamiento se congela antes de probarse en familias externas. En GPU, los kernels próximos al punto de inflexión se marcan como ambiguos sin excluirlos. El emparejamiento, las pruebas y los criterios de refinamiento se especifican en el Anexo \ref{anexo:metodologia}.

\section{Fase 3: agente de control DVFS}
\label{sec:metodologia-fase3}

La Fase 3 no vuelve a entrenar clasificadores ni vuelve a escoger frecuencias. Recibe como entradas congeladas tres artefactos de la Fase 2: el modelo ONNX de cada dispositivo, su umbral de abstención y la tabla de política. El modelo y el umbral proceden de la validación LOFO y de la auditoría de invarianza. La tabla de política procede, de forma independiente, de la comparación pareada de EDP entre niveles de frecuencia. Esta separación evita que el desempeño del clasificador se interprete como evidencia de que una frecuencia determinada mejora el EDP, o que la política de frecuencia determine la selección del modelo. La Fase 3 implementa detección, decisión, actuación, coordinación y restauración con esos artefactos sin modificarlos; la Fase 4 evalúa el conjunto ya congelado.

\subsection{Arquitectura}
\label{sec:metodologia-agente-arquitectura}

El agente son dos procesos independientes en C++, uno por dispositivo, coordinados por una señal de un solo bit (Figura \ref{fig:met-agente}). Se separaron porque los bucles operan a escalas distintas (ventanas de milisegundos en CPU, fases de segundos en GPU), porque un fallo en uno no impide que el otro restaure el dispositivo que le corresponde, y porque los experimentos de un solo dispositivo no cargan la biblioteca de gestión del otro. Ambos leen los modelos, umbrales y tabla de política congelados al inicio de la ejecución; el lanzador los traduce a las banderas de cada proceso.

\begin{figure}[H]
\centering
\begin{tikzpicture}[
    font=\footnotesize,
    b/.style={draw, rounded corners=2pt, align=center, minimum height=1.0cm, inner sep=3pt, text width=4.6cm},
    cpu/.style={b, fill=blue!7, draw=blue!45!black},
    gpu/.style={b, fill=orange!10, draw=orange!55!black},
    tab/.style={b, fill=green!10, draw=green!45!black, text width=9.0cm},
    fl/.style={-{Latex[length=2.2mm]}, semithick},
    lb/.style={font=\scriptsize, align=center},
]
    \node[tab] (t) {Tabla de política y modelos ONNX (Fase 2)};
    \node[cpu, below left=1.0cm and -1.6cm of t] (c1) {\textbf{Proceso de CPU}\\contadores por ventana\\XGBoost, histéresis};
    \node[gpu, below right=1.0cm and -1.6cm of t] (g1) {\textbf{Proceso de GPU}\\NVML, ventana sostenida\\random forest, abstención};
    \node[cpu, below=0.9cm of c1] (c2) {Actuador de CPU\\rango por núcleo y turbo};
    \node[gpu, below=0.9cm of g1] (g2) {Actuador de GPU\\bloqueo de reloj};
    \draw[fl] (t) -| (c1); \draw[fl] (t) -| (g1);
    \draw[fl] (c1) -- (c2); \draw[fl] (g1) -- (g2);
    \draw[fl, dashed] (g1.west) -- node[lb, above] {señal de\\GPU activa} (c1.east);
\end{tikzpicture}
\caption[Arquitectura del agente]{Arquitectura del agente: dos procesos, una tabla de política compartida y una señal de coordinación. Fuente: elaboración propia.}
\label{fig:met-agente}
\end{figure}

Cada proceso admite los tres brazos de la Tabla \ref{tab:agente-brazos}. El brazo \textit{base} es la referencia contra la que se compara todo, y el brazo \textit{sombra} mide la sobrecarga del agente sin su efecto sobre el hardware.

\begin{table}[H]
\centering
\caption{Brazos del agente.}
\label{tab:agente-brazos}
\small
\begin{tabular}{@{}>{\raggedright\arraybackslash}p{2.2cm}>{\raggedright\arraybackslash}p{3.3cm}>{\raggedright\arraybackslash}p{7.3cm}@{}}
\toprule
\textbf{Brazo} & \textbf{Agente} & \textbf{Comportamiento} \\
\midrule
Base & No corre & El dispositivo queda bajo su \textit{governor} nativo (\texttt{intel\_pstate} con turbo en CPU, gestión automática del controlador en GPU) \\
Sombra & Muestrea, clasifica y decide & Registra cada decisión y nunca escribe sobre el hardware \\
Activo & Muestrea, clasifica, decide y escribe & Aplica la frecuencia decidida y restaura el estado al terminar \\
\bottomrule
\end{tabular}
\end{table}

\subsection{Detección de fase, decisión y parámetros}
\label{sec:metodologia-agente-gpu}

El agente de GPU delimita fases con la utilización NVML, reúne 3~s de actividad y decide una vez por fase. Si se abstiene, libera el reloj. El agente de CPU decide por ventana y exige 50 propuestas consecutivas antes de cambiar de nivel. Los umbrales y periodos se reúnen en la Tabla \ref{tab:met-parametros}; la máquina de decisión se verificó por separado de la biblioteca de inferencia.

\begin{table}[H]
\centering
\caption{Parámetros del agente.}
\label{tab:met-parametros}
\small
\setlength{\tabcolsep}{3pt}
\begin{tabular}{@{}>{\raggedright\arraybackslash}p{1.4cm}>{\raggedright\arraybackslash}p{3.1cm}>{\raggedright\arraybackslash}p{8.0cm}@{}}
\toprule
\textbf{Proceso} & \textbf{Parámetro} & \textbf{Valor y razón} \\
\midrule
GPU & Umbral de actividad & 5\% de utilización \\
GPU & Ventana sostenida & 3~s; una decisión por fase \\
GPU & Cadencia de muestreo & 50~ms \\
GPU & Umbral de abstención & 0.90 (cobertura media por celda $\geq$ 0.60); seleccionado en la Fase 2 (sección \ref{sec:resultados-gpu-invarianza}) \\
GPU & Permanencia mínima entre cambios & 3.7~s, 10 veces la latencia de conmutación medida en la Fase 3 (sección \ref{sec:resultados-agente-latencias}) \\
CPU & Umbral de decisión & 0.85; congelado en la Fase 2 (sección \ref{sec:resultados-cpu-selectiva}) \\
CPU & Histéresis & 50 ventanas consecutivas \\
CPU & Piso con GPU activa & 3.2~GHz; verificado en la Fase 3 (sección \ref{sec:resultados-agente-restauracion}) \\
\bottomrule
\end{tabular}
\end{table}

\subsection{Política, actuación y coordinación}
\label{sec:metodologia-agente-politica}

La Tabla \ref{tab:agente-politica} separa la política operativa derivada en la Fase 2 de las variantes experimentales de CPU. La salida operativa de la Fase 2 para CPU es \textit{no actuar} en ambas clases, por lo que no se presenta F0 ni F1 como recomendación de despliegue. Para comprobar que esa decisión no oculta un efecto al cerrar el lazo de control, la Fase 4 incorpora un brazo de CPU experimental que fija F0 en cómputo y F1 en memoria, y una variante \textit{activo-F0} que mantiene F0 en ambas. Estas variantes no actualizan la tabla derivada ni alteran el modelo; se comparan contra el \textit{governor} nativo.

\begin{table}[H]
\centering
\caption{Acciones del agente. Las filas de GPU provienen de la política derivada en la Fase 2; las filas de CPU son variantes experimentales evaluadas en la Fase 4.}
\label{tab:agente-politica}
\small
\setlength{\tabcolsep}{4pt}
\begin{tabular}{@{}>{\raggedright\arraybackslash}p{1.7cm}>{\raggedright\arraybackslash}p{3.0cm}>{\raggedright\arraybackslash}p{7.5cm}@{}}
\toprule
\textbf{Dispositivo} & \textbf{Clase} & \textbf{Acción} \\
\midrule
CPU & \textit{compute\_bound} & Fija F0 (3.2~GHz, turbo desactivado); experimental \\
CPU & \textit{memory\_bound} & Baja a F1 (2.9~GHz) con confianza mínima de 0.85; experimental \\
GPU & \textit{memory\_bound} & Fija el reloj en 1260~MHz (F1) \\
GPU & \textit{compute\_bound} & No actúa: libera el reloj al \textit{governor} nativo \\
GPU & \textit{revisar} (abstención) & No actúa: libera el reloj al \textit{governor} nativo \\
\bottomrule
\end{tabular}
\end{table}

El actuador de CPU reutiliza las reglas de restricción de rango por núcleo del módulo de la Fase 1 (delegados y hermanos de \textit{SMT}), desactiva el turbo global con el envoltorio administrativo del nodo y escribe los núcleos en paralelo. Registra el estado previo, verifica cada escritura por relectura, falla cerrado ante cualquier inconsistencia y restaura el estado inicial al terminar, incluso ante \texttt{SIGTERM} o \texttt{SIGINT}; su equivalencia con el módulo de la Fase 1 se comprobó con una prueba de paridad. El actuador de GPU fija y libera el reloj con la herramienta de administración del controlador, con la misma disciplina de restauración.

La coordinación entre procesos es una regla de protección distinta de la tabla de política. Cuando la señal de actividad GPU está presente, impone temporalmente un piso de CPU de 3.2~GHz y no espera la histéresis. No reclasifica la fase de CPU ni convierte F0 en una acción recomendada por la Fase 2. Se adopta para evitar que las señales de una CPU que espera a la GPU desencadenen una reducción de frecuencia que prolongue la fase GPU. La Fase 3 verifica que el piso se aplique y se restaure; su efecto sobre EDP se mide por separado en la Fase 4 con el brazo sin piso.

\subsection{Aplicaciones compuestas y verificación previa}
\label{sec:metodologia-agente-compuestas}

Las aplicaciones compuestas registran con el reloj monotónico las fronteras reales de fase para puntuar al agente (Tabla \ref{tab:agente-compuestas}). A usa familias vistas y B familias inéditas; una tercera aplicación alterna CPU y GPU. Cada fase supera los 3~s de observación y se separa de la siguiente por 1~s de reposo. Los kernels inéditos de GPU usan parámetros ampliados del mismo binario verificado para sostener la actividad.

\begin{table}[H]
\centering
\caption{Aplicaciones compuestas con verdad de fase.}
\label{tab:agente-compuestas}
\small
\setlength{\tabcolsep}{4pt}
\begin{tabular}{@{}>{\raggedright\arraybackslash}p{1.4cm}>{\raggedright\arraybackslash}p{2.0cm}>{\raggedright\arraybackslash}p{10.2cm}@{}}
\toprule
\textbf{Alcance} & \textbf{Aplicación} & \textbf{Fases (clase)} \\
\midrule
CPU & A (vistos) & \texttt{dgemm\_n2048} (\textit{compute}), \texttt{npb\_cg} (\textit{memory}) \\
CPU & B (inéditos) & \texttt{cpu\_xsbench\_omp} (\textit{memory}), \texttt{cpu\_rsbench\_omp} (\textit{compute}) \\
GPU & A (vistos) & \texttt{gpu\_cutlass\_simt\_dgemm\_n4096} (\textit{compute}), \texttt{gpu\_rajaperf\_stream\_triad} (\textit{memory}) \\
GPU & B (inéditos) & BabelStream con \texttt{--numtimes 1000} (\textit{memory}), \texttt{rodinia\_lavamd} con \texttt{-boxes1d 100} (\textit{compute}) \\
\bottomrule
\end{tabular}
\end{table}

Antes de la evaluación se verificó el agente por dispositivo y con ambos a la vez: latencia de conmutación, sobrecarga, restauración bajo terminación anómala inducida y cumplimiento del piso de coordinación (sección \ref{sec:resultados-agente}).

\section{Fase 4: validación experimental}
\label{sec:metodologia-fase4}

\subsection{Diseño de la matriz}
\label{sec:metodologia-fase4-matriz}

La matriz inicial compara el agente contra el \textit{governor} nativo en tres ejes: el alcance (solo CPU, solo GPU o ambos a la vez sobre una aplicación que alterna fases), el tipo de kernel (vistos frente a inéditos, que mide la brecha de generalización con las mismas celdas) y el brazo (Tabla \ref{tab:fase4-matriz}). El orden de las celdas se aleatoriza con semilla fija y cada una ejecuta la aplicación dos ciclos completos. El brazo \textit{sombra} aísla la sobrecarga del agente; \textit{activo-F0} y \textit{activo sin piso} se incluyen como brazos experimentales porque la política de CPU resultó \textit{no actuar} en la Fase 2 y se quiere medir si esa conclusión se sostiene con el agente en el lazo. La matriz inicial usó cinco repeticiones en el alcance conjunto y tres en los demás casos. Con tres observaciones por brazo, el análisis es descriptivo mediante razones de medianas y valores individuales; una comparación no apareada de tres contra tres no puede producir evidencia exacta bilateral inferior a 0.10.

\begin{table}[H]
\centering
\caption{Matriz experimental de la Fase 4.}
\label{tab:fase4-matriz}
\small
\setlength{\tabcolsep}{3pt}
\begin{tabular}{@{}>{\raggedright\arraybackslash}p{1.7cm}>{\raggedright\arraybackslash}p{6.3cm}>{\centering\arraybackslash}p{1.1cm}>{\centering\arraybackslash}p{1.3cm}>{\centering\arraybackslash}p{2.1cm}@{}}
\toprule
\textbf{Alcance} & \textbf{Brazos} & \textbf{Vistos} & \textbf{Inéditos} & \textbf{Repeticiones} \\
\midrule
CPU & base, sombra, activo, activo-F0 & A & B & 3 \\
GPU & base, sombra, activo & A & B & 3 \\
Conjunto & base, sombra, activo con piso, activo sin piso & A & B & 3 (5 en la matriz inicial) \\
\bottomrule
\end{tabular}
\end{table}

\subsection{Medición y análisis}
\label{sec:metodologia-fase4-medicion}

En todos los escenarios, cada celda registra duración, energía de CPU (RAPL de ambos paquetes) y energía de GPU (contador NVML) por separado. Así, el EDP del nodo, $\mathrm{EDP} = (E_{\mathrm{CPU}} + E_{\mathrm{GPU}})\,T$, se calcula y se desagrega después; también se calcula $E_{\mathrm{GPU}}\,T$. La ventana energética abarca desde el arranque de la aplicación hasta su fin, con los procesos del agente ya activos, por lo que incluye su sobrecarga. Cada celda verifica la finalización de la aplicación y de los procesos del agente que correspondan, y restaura los estados de frecuencia que modificó.

En la matriz y en el escenario E, las aplicaciones compuestas registran las fronteras de fase conocidas y las decisiones de cada proceso (instante, clase, confianza y escritura); la aplicación se fija a los núcleos 0 a 3 y se libera el reloj de GPU entre celdas. LAMMPS también se ejecuta en esos núcleos, pero sus fases son naturales y no se registran fronteras. CloverLeaf se ejecuta como un proceso CUDA Fortran continuo con ocho CPU asignadas; no se le impone ni se le atribuye una frontera de fase. Para CloverLeaf se conserva además la salida numérica, muestras NVML cada 250~ms y el estado de reloj antes y después de cada celda.

Para cada combinación de alcance, kernels y brazo se calcula la mediana de duración, energías y EDP y su razón frente al brazo \textit{base} (menor es mejor). En la matriz inicial conjunta, que tiene cinco observaciones no apareadas por brazo, el EDP se contrasta con Mann--Whitney bilateral. Los dos confirmatorios de cinco bloques, E-A y CloverLeaf, se analizan por bloque: se informan las razones por bloque y la prueba exacta bilateral de signos. Con cinco pares, incluso cinco diferencias del mismo signo dan $p=0.0625$, por lo que el valor se reporta sin declarar significancia al nivel 0.05.

La clasificación se puntúa contra las fronteras reales solo en las aplicaciones compuestas: cada decisión se atribuye a la fase que contiene su instante y se reportan correctas, erróneas y abstenciones por alcance, tipo de kernel y dispositivo, junto con la cobertura. La diferencia de exactitud entre A y B mide la generalización dentro de ese protocolo. En CPU las decisiones son por ventana, así que la exactitud se pondera por ventanas y no por fases. En LAMMPS y CloverLeaf se reportan las decisiones y la actuación observada, pero no una exactitud de clasificación, pues no existe una verdad de fase independiente con esa granularidad.

\subsection{Escenarios adicionales}
\label{sec:metodologia-fase4-escenarios}

Además de la matriz inicial, que se reporta completa y es la referencia, se ejecutaron los protocolos complementarios de la Tabla \ref{tab:fase4-mapa}; sus condiciones completas se conservan en la Tabla \ref{tab:fase4-escenarios} del Anexo \ref{anexo:metodologia}. Cada protocolo se fijó antes de sus mediciones y todos los ejecutados se reportan, cualquiera que sea su resultado. El escenario E se definió con el criterio de selección de sus familias (ganancia de EDP con F1 de al menos 10\% en la Fase 2, más una variante con familias no vistas) y se verificó con una corrida previa de una repetición por brazo antes de la matriz completa.

La referencia secundaria sin turbo se incorporó en las campañas complementarias C y E. En ese brazo no se inicia ningún daemon: se desactiva el turbo de CPU inmediatamente antes de la ventana de medición, se verifica el estado solicitado y se restaura al terminar la celda. Se ejecutan tres repeticiones para cada combinación de aplicación y alcance incluida. El control conserva el \textit{governor} nativo y el reloj nativo de la GPU, por lo que solo permite atribuir al turbo de CPU una diferencia entre referencias; no representa una frecuencia fija de GPU.

La Tabla \ref{tab:met-amenazas} resume las principales amenazas a la validez y la medida que la controla.

\begin{table}[H]
\centering
\caption{Mapa de los escenarios de la Fase 4. El protocolo completo se conserva en el Anexo \ref{anexo:metodologia}.}
\label{tab:fase4-mapa}
\small
\begin{tabular}{@{}p{3.1cm}p{10.3cm}@{}}
\toprule
\textbf{Escenario} & \textbf{Contraste principal} \\
\midrule
A & Matriz inicial por dispositivo, familias vistas e inéditas y brazo \\
C & Agente de CPU con un hilo de inferencia \\
E & GPU con fases largas de memoria, familias vistas (E-A) e inéditas (E-B) \\
REF sin turbo & Línea base secundaria para aislar el turbo de CPU \\
D & LAMMPS con cuatro entradas y fases naturales \\
F & Premedición diagnóstica de la solicitud de F1 fijo \\
Confirmatorio E-A & Cinco bloques con REF, sombra, agente y solicitud de F1 fijo \\
CloverLeaf & Cinco bloques con REF, sombra y agente en aplicación continua \\
\bottomrule
\end{tabular}
\end{table}
